\subsection{两个测度的乘积}

定理 1. --- 设 X 与 Y 是两个局部紧空间， \(\lambda\) 是 X 上的一个测度， \(\mu\) 是 Y 上的一个测度；则存在唯一的测度 \(\nu\) （X \(\times\) Y 上的），使得对每个函数 \(g \in \mathcal{K}(\mathrm{X};\mathbf{C})\) 及每个函数 \(h \in \mathcal{K}(\mathrm{Y};\mathbf{C})\) ，

\[
\int g (x) h (y) d \nu (x, y) = \left(\int g (x) d \lambda (x)\right) \left(\int h (y) d \mu (y)\right).\tag{1}
\]

引理 1. --- 设 X 与 Y 是两个局部紧空间，K（相应地 L）是 X（相应地 Y）的一个紧子集。

\begin{enumerate}
\def\labelenumi{(\roman{enumi})}
\tightlist
\item
  典范双射在 \(\mathcal{K}(\mathrm{X}\times \mathrm{Y},\mathrm{K}\times \mathrm{L};\mathbf{C})\) 上的限制
\end{enumerate}

\[
\omega : \mathcal {F} (\mathrm{X} \times \mathrm{Y}; \mathbf {C}) \rightarrow \mathcal {F} \bigl (\mathrm{X}; \mathcal {F} (\mathrm{Y}; \mathbf {C}) \bigr)
\]

（S, R, §4, No.~14）是 Banach 空间 \(\mathcal{K}(\mathrm{X}\times\mathrm{Y},\mathrm{K}\times\mathrm{L};\mathbf{C})\) 到 Banach 空间 \(\mathcal{K}(\mathrm{X},\mathrm{K};\mathcal{K}(\mathrm{Y},\mathrm{L};\mathbf{C}))\) 上的一个等距。

\begin{enumerate}
\def\labelenumi{(\roman{enumi})}
\setcounter{enumi}{1}
\tightlist
\item
  向量空间 \(\mathcal{K}(\mathrm{X},\mathrm{K};\mathbf{C})\otimes_{\mathbf{C}}\mathcal{K}(\mathrm{Y},\mathrm{L};\mathbf{C})\) （典范地等同于 \(\mathcal{K}(\mathrm{X}\times \mathrm{Y},\mathrm{K}\times \mathrm{L};\mathbf{C})\) 的一个子空间；A, II, §7, No.~7, 命题 15 的推论之后的注记）在这个 Banach 空间中稠密。
\end{enumerate}

\(\omega\) 下 \(\mathcal{K}(X\times Y,K\times L;C)\) 的像含于 \(\mathcal{K}(X,K;\mathcal{K}(Y,L;C))\) ，这是直接的。反之，若 u 是 X 到 \(\mathcal{K}(Y,L;C)\) 中的一个连续映射，其支集含于 K，则映射 \((x,y)\mapsto\mathbf{u}(x)(y)\) （ \(X\times Y\) 到 C 中的）是连续的且支集含于 \(K\times L\) ，因此 \(\omega\) 在 \(\mathcal{K}(X\times Y,K\times L;C)\) 上的限制是此空间到 \(\mathcal{K}(X,K;\mathcal{K}(Y,L;C))\) 上的一个双射；这一限制是 Banach 空间等距这一事实由关系式

\[
\sup _ {(x, y) \in \mathrm{K} \times \mathrm{L}} | f (x, y) | = \sup _ {x \in \mathrm{K}} \left(\sup _ {y \in \mathrm{L}} | f (x, y) |\right).
\]

得出。这就证明了 (i)；另一方面，在 \(\omega\) 下，

\[
\mathcal {K} (\mathrm{X}, \mathrm{K}; \mathbf {C}) \otimes_ {\mathbf {C}} \mathcal {K} (\mathrm{Y}, \mathrm{L}; \mathbf {C})
\]

（等同于 \(\mathcal{K}(\mathrm{X}\times \mathrm{Y},\mathrm{K}\times \mathrm{L};\mathbf{C})\) 的一个子空间）的像仍是空间 \(\mathcal{K}(\mathrm{X},\mathrm{K};\mathbf{C})\otimes_{\mathbf{C}}\mathcal{K}(\mathrm{Y},\mathrm{L};\mathbf{C})\) ，不过这次是典范地等同于 X 到 \(\mathcal{K}(\mathrm{Y},\mathrm{L};\mathbf{C})\) 中的映射所成的一个空间（A, II, §7, No.~7, 命题 15 的推论）；而 \(\mathcal{K}(\mathrm{X},\mathrm{K};\mathcal{K}(\mathrm{Y},\mathrm{L};\mathbf{C}))\) 的这一子空间已知在后一空间中稠密（§1, No.~2, 命题 5），于是 (ii) 的结论由 \(\omega\) 的这一限制是拓扑同构这一事实得出。

引理证毕之后，我们现在注意， \(X \times Y\) 的每个紧子集都含于某个乘积 \(K \times L\) 之中，其中 K（相应地 L）是 X（相应地 Y）的一个紧子集。于是由引理 1, (ii) 可知，子空间 \(\mathcal{K}(X; \mathbf{C}) \otimes_{\mathbf{C}} \mathcal{K}(Y; \mathbf{C})\) 在 \(\mathcal{K}(X \times Y; \mathbf{C})\) 中稠密；由于关系式 (1) 也可写成 \(\nu(g \otimes h) = \lambda(g)\mu(h)\) （对 \(g \in \mathcal{K}(X; \mathbf{C})\) ， \(h \in \mathcal{K}(Y; \mathbf{C})\) ），故 \(\nu\) 的唯一性立即得出。为证 \(\nu\) 的存在性，我们将使用下述引理：

引理 2. --- 在引理 1 的记号下，对每个函数 \(f \in \mathcal{K}(\mathrm{X} \times \mathrm{Y}, \mathrm{K} \times \mathrm{L}; \mathbf{C})\) ，函数

\[
y \mapsto h (y) = \int f (x, y) d \lambda (x)\tag{2}
\]

属于 \(\mathcal{K}(\mathrm{Y},\mathrm{L};\mathbf{C})\)

对每个函数 \(\mathbf{u} \in \mathcal{K}(\mathrm{X};\mathcal{K}(\mathrm{Y},\mathrm{L};\mathbf{C}))\) ，积分 \(\int \mathbf{u}(x)d\lambda (x)\) 属于 \(\mathcal{K}(\mathrm{Y},\mathrm{L};\mathbf{C})\) ，因为后者是 Banach 空间（ \(\S 3\) , No.~3, 命题 7 的推论 1）；而对 \(\mathbf{u} = \omega(f)\) 及每个 \(y \in \mathrm{Y}\) ，

\[
\left\langle \int \mathbf {u} (x) d \lambda (x), \varepsilon_ {y} \right\rangle = \int \mathbf {u} (x) (y) d \lambda (x) = \int f (x, y) d \lambda (x),
\]

由此得该引理。

于是，对每个函数 \(f \in \mathcal{K}(\mathrm{X} \times \mathrm{Y}; \mathbf{C})\) ，考虑数 \(\nu(f) = \mu\left(\int f(x, y) d\lambda(x)\right)\) （出于记号的滥用，我们也把它写作 \(\int d\mu(y) \int f(x, y) d\lambda(x)\) ）；若 K（相应地 L）是 X（相应地 Y）的一个紧子集，则存在一个数 \(a_{K}\) （相应地 \(b_{L}\) ），使得对每个函数 \(g \in \mathcal{K}(\mathrm{X}, \mathrm{K}; \mathbf{C})\) （相应地 \(h \in \mathcal{K}(\mathrm{Y}, \mathrm{L}; \mathbf{C})\) ）有 \(|\lambda(g)| \leqslant a_{\mathrm{K}} \|g\|\) （相应地 \(|\mu(h)| \leqslant b_{\mathrm{L}} \|h\|\) ）。由此可知，对每个函数 \(f \in \mathcal{K}(\mathrm{X} \times \mathrm{Y}, \mathrm{K} \times \mathrm{L}; \mathbf{C})\) ，

\[
\left| \int f (x, y) d \lambda (x) \right| \leqslant a _ {K} \| f \|
\]

对每个 \(y \in Y\) 成立，由此得 \(|\nu(f)| \leqslant a_{K}b_{L}\|f\|\) 。于是线性形式 \(\nu\) （ \(\mathcal{K}(X \times Y; \mathbf{C})\) 上的）是 \(X \times Y\) 上的一个测度，它显然满足 (1)，这就完成了 Th. 1 的证明。

定义 1. --- 给定分别定义在两个局部紧空间 X、Y 上的两个测度 \(\lambda\) 、 \(\mu\) ，我们把 \(\lambda\) 与 \(\mu\) 的乘积测度定义为唯一的测度 \(\nu\) （ \(X \times Y\) 上的），它对每个函数 \(g \in \mathcal{K}(X; \mathbf{C})\) 及每个函数 \(h \in \mathcal{K}(Y; \mathbf{C})\) 满足关系式 (1)。

在 Th. 1 的证明中，X 与 Y 两空间的地位可以互换；把 \(Y \times X\) 与 \(X \times Y\) 典范地等同起来，我们便在 \(X \times Y\) 上定义了测度

\[
f \mapsto \int d \lambda (x) \int f (x, y) d \mu (y)
\]

它仍满足条件 (1)。于是我们证明了下述定理：

定理 2. --- 设 \(\lambda\) 、 \(\mu\) 是分别定义在两个局部紧空间 \(X\) 、 \(Y\) 上的两个测度。对每个函数 \(f\) （ \(\mathcal{K}(X \times Y; \mathbf{C})\) 中的），\(f\) 关于乘积测度 \(\nu\) （\(\lambda\) 与 \(\mu\) 的乘积）的积分的值为

\[
\begin{array}{c} \int f (x, y) d \nu (x, y) = \int d \lambda (x) \int f (x, y) d \mu (y) \\ = \int d \mu (y) \int f (x, y) d \lambda (x). \end{array}\tag{3}
\]

由于最后的公式，f 关于乘积测度 \(\nu\) 的积分通常记作 \(\iint f d\lambda d\mu\) ，或 \(\iint f d\mu d\lambda\) ，或 \(\iint f \lambda \mu\) ，或 \(\iint f \mu \lambda\) ，或 \(\iint f(x, y) d\lambda(x) d\mu(y)\) ，或 \(\iint f(x, y) d\mu(y) d\lambda(x)\) ，或 \(\iint f(x, y) \lambda(x)\mu(y)\) ，或 \(\iint f(x, y) \mu(y)\lambda(x)\) ；称它为 f 关于 \(\lambda\) 与 \(\mu\) 的二重积分。采用这一记号，公式 (3) 可写成

\[
\begin{array}{c} \iint f (x, y) d \lambda (x) d \mu (y) = \int d \lambda (x) \int f (x, y) d \mu (y) \\ = \int d \mu (y) \int f (x, y) d \lambda (x). \end{array}\tag{4}
\]

公式 (3) 表明，若 \(\lambda\) 与 \(\mu\) 是实（相应地，正）测度，则乘积测度 \(\nu\) 是实（相应地，正）测度。

例子. --- 1) Dirac 测度 \(\varepsilon_{x}\) （ \(x \in X\) ）与 \(\varepsilon_{y}\) （ \(y \in Y\) ）的乘积测度是 Dirac 测度 \(\varepsilon_{(x,y)}\) 。

\begin{enumerate}
\def\labelenumi{\arabic{enumi})}
\setcounter{enumi}{1}
\tightlist
\item
  取 \(\mathbf{X} = \mathbf{Y} = \mathbf{R}\) ，并取 \(\lambda\) 与 \(\mu\) 为 \(\mathbf{R}\) 上的 Lebesgue 测度（§1, No.~3）；其乘积称为 \(\mathbf{R}^2\) 上的 Lebesgue 测度；函数 \(f \in \mathcal{K}(\mathbf{R}^2; \mathbf{C})\) 关于该测度的积分记作 \(\iint f(x, y) dx dy\) 或 \(\iint f(x, y) dy dx\) ；公式 (4) 对在紧矩形 \([a, b] \times [c, d]\) 之外为零的函数给出公式
\end{enumerate}

\[
\int_ {c} ^ {d} d y \int_ {a} ^ {b} f (x, y) d x = \int_ {a} ^ {b} d x \int_ {c} ^ {d} f (x, y) d y
\]

这已在 FRV, II, §3, No.~6 中证明。

由于 \(\mathbf{R}\) 上的 Lebesgue 测度在每个平移下不变（§1, No.~3），立即得出 \(\mathbf{R}^2\) 上的 Lebesgue 测度在 \(\mathbf{R}^2\) 的每个平移下不变。

注. --- 设 E 是一个 Hausdorff 局部凸空间，f 是 \(\mathcal{K}(X \times Y; E)\) 中的一个映射，使得 \(\mathbf{f}(X \times Y)\) 含于 E 的一个完备凸子集 C。于是对每个 \(y \in Y\) ，积分 \(\mathbf{h}(y) = \int \mathbf{f}(x, y) d\lambda(x)\) 属于 E（§3, No.~3, 命题 7）；此外，函数 h 属于 \(\widetilde{\mathcal{K}}(Y; E)\) ：因为对每个 \(z' \in E'\) 有

\[
\langle \mathbf {h} (y), \mathbf {z} ^ {\prime} \rangle = \int \langle \mathbf {f} (x, y), \mathbf {z} ^ {\prime} \rangle d \lambda (x),
\]

因而由引理 2， \(y \mapsto \langle \mathbf{h}(y), \mathbf{z}' \rangle\) 属于 \(\mathcal{K}(\mathrm{Y}; \mathbf{C})\) 。于是积分 \(\int h d\mu\) 有定义（且先验地属于 \(E'^{*}\) ）；我们证明

\[
\begin{array}{r l} \iint \mathbf {f} (x, y) d \lambda (x) d \mu (y) & = \int d \mu (y) \int \mathbf {f} (x, y) d \lambda (x) \\ & = \int d \lambda (x) \int \mathbf {f} (x, y) d \mu (y), \end{array}\tag{5}
\]

从而把公式 (4) 加以推广。事实上，对每个 \(\mathbf{z}' \in \mathrm{E}'\) 有

\[
\begin{array}{r l} \left\langle \iint \mathbf {f} d \lambda d \mu , \mathbf {z} ^ {\prime} \right\rangle & = \iint \langle \mathbf {f}, \mathbf {z} ^ {\prime} \rangle d \lambda d \mu = \int d \mu \int \langle \mathbf {f}, \mathbf {z} ^ {\prime} \rangle d \lambda \\ & = \int \left\langle \int \mathbf {f} d \lambda , \mathbf {z} ^ {\prime} \right\rangle d \mu = \left\langle \int d \mu \int \mathbf {f} d \lambda , \mathbf {z} ^ {\prime} \right\rangle \end{array}
\]

这是依据 (4)，由此得 (5)。

